\documentclass[11pt]{article}
\usepackage[a4paper,margin=2.3cm]{geometry}
\usepackage{amsmath,amssymb}
\usepackage{fancyhdr}
\usepackage[most]{tcolorbox}
\usepackage{enumitem}
\usepackage{xcolor}
\usepackage{parskip}

\definecolor{introblue}{RGB}{230,240,250}
\definecolor{introbluebar}{RGB}{70,130,180}
\definecolor{curiousyellow}{RGB}{255,248,225}
\definecolor{curiousbar}{RGB}{200,160,40}
\definecolor{sanitygreen}{RGB}{232,245,233}
\definecolor{sanitybar}{RGB}{60,140,90}
\definecolor{revisiongray}{RGB}{240,240,240}
\definecolor{revisionbar}{RGB}{120,120,120}

\newtcolorbox{intuicao}{
  colback=introblue, colframe=introbluebar, colbacktitle=introbluebar, coltitle=white,
  boxrule=0pt, leftrule=3pt, arc=1pt, left=8pt, right=8pt, top=5pt, bottom=5pt,
  fonttitle=\bfseries, title={INTUIÇÃO}
}
\newtcolorbox{curiosos}{
  colback=curiousyellow, colframe=curiousbar, colbacktitle=curiousbar, coltitle=white,
  boxrule=0pt, leftrule=3pt, arc=1pt, left=8pt, right=8pt, top=5pt, bottom=5pt,
  fonttitle=\bfseries, title={PARA OS CURIOSOS --- de onde vem essa fórmula}
}
\newtcolorbox{sanidade}{
  colback=sanitygreen, colframe=sanitybar, colbacktitle=sanitybar, coltitle=white,
  boxrule=0pt, leftrule=3pt, arc=1pt, left=8pt, right=8pt, top=5pt, bottom=5pt,
  fonttitle=\bfseries, title={CHECAGEM DE SANIDADE}
}
\newtcolorbox{revisao}{
  colback=revisiongray, colframe=revisionbar, colbacktitle=revisionbar, coltitle=white,
  boxrule=0pt, leftrule=3pt, arc=1pt, left=8pt, right=8pt, top=5pt, bottom=5pt,
  fonttitle=\bfseries\itshape, title={NOTA DE REVISÃO}
}

\pagestyle{fancy}
\fancyhf{}
\fancyhead[L]{\small Do score ao erro}
\fancyhead[R]{\small Prof. Thiago M. Paixão}
\renewcommand{\headrulewidth}{0.3pt}
\fancyfoot[C]{\small\thepage}

\newcommand{\R}{\mathbb{R}}

\title{\vspace{-1.5cm}\Large\bfseries Do Score ao Erro\\[4pt]\normalsize\normalfont Softmax, Perda e Gradientes --- um guia intuitivo}
\author{Prof. Thiago M. Paixão \\ \small Ifes -- Campus Serra}
\date{}

\begin{document}
\maketitle
\vspace{-1cm}

Pense num sistema de ranking qualquer --- um placar de partida, um retorno de API que atribui uma pontuação a cada opção. É basicamente isso que uma rede neural de classificação faz na última camada: para cada classe possível, ela calcula um score (chamado de \emph{logit}). Só que scores brutos não são probabilidades --- podem ser negativos, não somam 1, não são comparáveis entre si de forma direta. Esta nota constrói, passo a passo, a tubulação que transforma esses scores em probabilidades e, a partir delas, em um sinal de erro que a rede usa para aprender.

Vamos do caso mais simples (um único exemplo) até a forma que você de fato escreve em código (mini-batches vetorizados) --- e, no caminho, vamos abrir o capô e calcular os gradientes manualmente, sem atalhos, para que a mágica do \texttt{loss.backward()} pare de ser mágica.

\section{Um único exemplo}

Seja $x \in \R^d$ um vetor de entrada (por exemplo, os pixels de uma imagem já achatados em um vetor) e $C$ o número de classes. Cada classe $k$ tem seu próprio detector de score: um vetor de pesos $w_k \in \R^d$ e um viés escalar $b_k$. Salvo indicação em contrário, todo vetor nesta nota é um \textbf{vetor coluna} --- isso inclui $x$ e cada $w_k$.

$$
W = \begin{bmatrix} w_1^\top \\ \vdots \\ w_C^\top \end{bmatrix} \in \R^{C\times d},
\qquad
b = \begin{bmatrix} b_1 \\ \vdots \\ b_C \end{bmatrix} \in \R^C
$$

\begin{revisao}
Como $w_k$ é vetor coluna ($d\times 1$), empilhar os $w_k$ como \emph{linhas} de $W$ exige transpô-los: a linha $k$ de $W$ é $w_k^\top$ ($1\times d$), não $w_k$. É essa transposição que faz $W \in \R^{C\times d}$ fechar dimensionalmente e que torna $z_k = w_k^\top x + b_k$ (Passo 1 abaixo) consistente com $z = Wx + b$.
\end{revisao}

\subsection*{Passo 1 --- Scores (logits)}
$$
z_k = w_k^\top x + b_k, \qquad k = 1,\dots,C
$$

Cada $z_k$ mede o quanto a entrada $x$ se parece com a classe $k$, segundo os pesos atuais. É um número real qualquer --- pode ser $-8.3$, pode ser $52.1$. Ainda não é probabilidade.

\begin{intuicao}
Pense em $w_k$ como um molde aprendido para a classe $k$. O produto $w_k^\top x$ é literalmente uma medida de similaridade (produto interno) entre a entrada e esse molde: quanto mais $x$ se alinha com $w_k$, maior o score $z_k$.
\end{intuicao}

\subsection*{Passo 2 --- Softmax: de scores a probabilidades}
$$
p_k = \mathrm{softmax}(z)_k = \frac{e^{z_k}}{\sum_{j=1}^{C} e^{z_j}}
$$

O softmax faz três coisas ao mesmo tempo: (i) exponencia cada score, tornando tudo positivo; (ii) normaliza pela soma, garantindo $\sum_k p_k = 1$; (iii) exagera as diferenças relativas --- um score um pouco maior que os outros vira uma probabilidade proporcionalmente muito maior, graças à exponencial.

\begin{intuicao}
É como transformar os pontos de um campeonato em probabilidade de título. Não basta ter mais pontos que o adversário: o quanto a mais importa, e o softmax captura exatamente essa não linearidade.
\end{intuicao}

\subsection*{Passo 3 --- A perda: penalizando a probabilidade errada}

Se $y$ é a classe correta do exemplo, a perda de entropia cruzada é:
$$
L = -\log(p_y)
$$

\begin{intuicao}
$-\log(\cdot)$ é uma função que vale $0$ quando o argumento é $1$, e explode para $+\infty$ quando o argumento tende a $0$. Ou seja: se o modelo acerta com $p_y \approx 1$, a perda é quase nula (sem castigo). Se o modelo erra feio, com $p_y \approx 0$ (achava que a classe correta era quase impossível), a perda dispara. É um professor rigoroso: só reclama quando você erra, e reclama muito mais quanto mais confiante você estava no erro.
\end{intuicao}

$$
x \;\longrightarrow\; z = Wx+b \;\longrightarrow\; p = \mathrm{softmax}(z) \;\longrightarrow\; L = -\log p_y
$$

Note que a perda depende de todos os $z_j$ (via o denominador do softmax), não apenas de $z_y$ --- isso vai importar na hora de derivar.

\section{Do exemplo para o dataset inteiro}

Com $N$ exemplos $\{(x^{(i)}, y^{(i)})\}_{i=1}^{N}$, cada um passa pela mesma tubulação (mesmos $W$, $b$, compartilhados), gerando seu próprio $z^{(i)}$, $p^{(i)}$ e perda individual $-\log p_{y^{(i)}}^{(i)}$.

A perda do dataset é simplesmente a média das perdas individuais:
$$
L(W,b) = \frac{1}{N}\sum_{i=1}^{N} -\log p^{(i)}_{y^{(i)}}
$$

\begin{intuicao}
$W$ e $b$ são os únicos botões que o treinamento pode girar --- e eles são compartilhados por todos os $N$ exemplos. Treinar é encontrar os valores de $W, b$ que, em média, deixam o time (todos os exemplos) satisfeito, não apenas um exemplo específico.
\end{intuicao}

\section{Do dataset para o mini-batch}

Calcular $L$ sobre os $N$ exemplos a cada passo de treino é caro (imagine $N=1{,}2$ milhão de imagens). Na prática, sorteia-se um subconjunto pequeno $\mathcal{B} \subset \{1,\dots,N\}$, de tamanho $B=|\mathcal{B}|$ (tipicamente 32, 64, 128\ldots), a cada iteração:
$$
L_{\mathcal{B}}(W,b) = \frac{1}{B}\sum_{i \in \mathcal{B}} -\log p^{(i)}_{y^{(i)}}
$$

\begin{intuicao}
$L_{\mathcal{B}}$ é uma amostra aleatória de $L$ --- como perguntar a $B$ pessoas escolhidas ao acaso, em vez da população inteira, para estimar uma opinião média. Não é exato, mas é um estimador não enviesado ($\mathbb{E}[L_{\mathcal{B}}] = L$) e muitíssimo mais barato de calcular. É essa troca --- precisão por velocidade --- que viabiliza o treinamento de redes grandes (SGD, Adam, etc.).
\end{intuicao}

\section*{Parte II --- Abrindo o capô: os gradientes}

Agora a pergunta que interessa para o treino: em que direção devo mexer cada peso para diminuir a perda? Essa direção é o gradiente. Vamos calculá-lo primeiro na mão, coordenada a coordenada --- e só depois compactar tudo em notação matricial.

\section{Gradiente --- versão não vetorizada}

Queremos $\dfrac{\partial L}{\partial w_k}$ e $\dfrac{\partial L}{\partial b_k}$, para cada classe $k=1,\dots,C$, ainda para um único exemplo.

\subsection*{Passo 1 --- gradiente em relação ao score $z_k$}

Aplicando a regra da cadeia através do softmax e do $-\log$, chega-se a um resultado surpreendentemente limpo:
$$
\frac{\partial L}{\partial z_k} = p_k - \mathbf{1}[k=y]
$$
onde $\mathbf{1}[k=y]$ vale 1 se $k$ é a classe correta, e 0 caso contrário.

\begin{curiosos}
O truque é reescrever $L$ de um jeito que ``abra'' o softmax em vez de escondê-lo dentro de uma fração. Denotando $S = \sum_{j=1}^{C} e^{z_j}$ (a soma no denominador do softmax, de modo que $p_j = e^{z_j}/S$):
$$
L = -\log(p_y) = -\log\!\left(\frac{e^{z_y}}{S}\right) = -\log(e^{z_y}) + \log(S) = -z_y + \log(S)
$$
(usamos a propriedade $\log(a/b) = \log a - \log b$, e $\log(e^{z_y}) = z_y$). Agora $L$ é uma soma de dois termos bem mais simples de derivar: um termo linear ($-z_y$) e um termo que só depende de $S$.

\textbf{Derivando o primeiro termo, $-z_y$, em relação a $z_k$:}

$z_y$ é apenas uma das coordenadas do vetor $z$ --- a da classe correta. Derivar $z_y$ em relação a $z_k$ só dá algo diferente de zero quando $k$ é exatamente essa coordenada:
$$
\frac{\partial(-z_y)}{\partial z_k} = -\mathbf{1}[k=y]
$$

\textbf{Derivando o segundo termo, $\log(S)$, em relação a $z_k$:}

Aqui usamos a regra da cadeia ($\frac{d}{du}\log(u) = \frac{1}{u}\cdot\frac{du}{dz_k}$) com $u=S$:
$$
\frac{\partial \log(S)}{\partial z_k} = \frac{1}{S}\cdot\frac{\partial S}{\partial z_k}
$$
Como $S = \sum_{j=1}^{C} e^{z_j}$, e apenas o termo $j=k$ dessa soma depende de $z_k$:
$$
\frac{\partial S}{\partial z_k} = \frac{\partial}{\partial z_k}\sum_{j=1}^{C} e^{z_j} = e^{z_k}
$$
Substituindo de volta:
$$
\frac{\partial \log(S)}{\partial z_k} = \frac{e^{z_k}}{S} = p_k
$$
(a última igualdade é só a própria definição de softmax --- $p_k = e^{z_k}/S$.)

\textbf{Somando os dois termos:}
$$
\frac{\partial L}{\partial z_k} = \underbrace{-\mathbf{1}[k=y]}_{\text{do termo } -z_y} + \underbrace{p_k}_{\text{do termo } \log(S)} = p_k - \mathbf{1}[k=y]
$$

Repare que não precisamos separar na mão os casos $k=y$ e $k\neq y$ --- o indicador $\mathbf{1}[k=y]$ já nasce naturalmente da derivada de $z_y$ e carrega essa distinção sozinho. Mas, para conferir, vale substituir os dois casos:
\begin{itemize}[leftmargin=1.3cm]
  \item Se $k=y$: $\partial L/\partial z_k = p_y - 1$ (negativo, pois $p_y \le 1$).
  \item Se $k\neq y$: $\partial L/\partial z_k = p_k - 0 = p_k$ (positivo, pois $p_k \ge 0$).
\end{itemize}
exatamente os dois comportamentos discutidos na caixa de intuição acima.
\end{curiosos}

\begin{intuicao}
Esse termo é literalmente o erro de cada score: a diferença entre o que o modelo previu ($p_k$) e o que deveria ter previsto (1 para a classe certa, 0 para as demais). Se o modelo acerta perfeitamente ($p_y=1$, $p_{k\neq y}=0$), esse erro é zero em todas as coordenadas --- e o gradiente desaparece, sinalizando nada a corrigir aqui.
\end{intuicao}

\subsection*{Passo 2 --- propagando até $w_k$ e $b_k$}

Como $z_k = w_k^\top x + b_k$, temos $\dfrac{\partial z_k}{\partial w_k} = x$ e $\dfrac{\partial z_k}{\partial b_k} = 1$. Pela regra da cadeia:
$$
\frac{\partial L}{\partial w_k} = \big(p_k - \mathbf{1}[k=y]\big)\, x, \qquad
\frac{\partial L}{\partial b_k} = p_k - \mathbf{1}[k=y]
$$

\begin{curiosos}
Formalmente, $L$ depende de $w_k$ através do vetor $z$ inteiro (não só de $z_k$), então a regra da cadeia multivariada exige, a rigor, um somatório sobre todas as coordenadas de $z$:
$$
\frac{\partial L}{\partial w_k} = \sum_{j=1}^{C} \frac{\partial L}{\partial z_j}\cdot\frac{\partial z_j}{\partial w_k}
$$
Mas cada classe tem seu próprio vetor de pesos, sem compartilhamento: $z_j = w_j^\top x + b_j$ depende apenas de $w_j$, nunca de $w_k$ com $k\neq j$. Logo
$$
\frac{\partial z_j}{\partial w_k} = \begin{cases} x, & j=k \\ 0, & j\neq k \end{cases}
$$
e o somatório colapsa num único termo sobrevivente ($j=k$):
$$
\frac{\partial L}{\partial w_k} = \frac{\partial L}{\partial z_k}\cdot x = \big(p_k - \mathbf{1}[k=y]\big)\,x
$$
É por isso que a passagem acima pôde ir direto de $\partial z_k/\partial w_k = x$ ao resultado final, sem mencionar os demais $C-1$ termos: eles são todos nulos.
\end{curiosos}

\begin{intuicao}
Se $k$ é a classe correta ($k=y$): o gradiente é $(p_y - 1)x$, e como $p_y \le 1$ esse fator é $\le 0$. O passo de descida do gradiente (que subtrai o gradiente) vai então aumentar $w_y^\top x$ --- empurrando $p_y$ para mais perto de 1. Faz sentido: queremos reforçar a classe certa.

Se $k$ é uma classe errada ($k\neq y$): o gradiente é $p_k \cdot x$, positivo. O passo de descida vai diminuir $w_k^\top x$ --- suprimindo a confiança nas classes erradas. O modelo está literalmente roubando probabilidade das classes erradas para dar à classe certa.
\end{intuicao}

\section{Forma vetorizada}

Empilhar $C$ equações escalares repetidas é sinal de que dá para escrever em forma vetorial. Defina o vetor de erro $\delta \in \R^C$:
$$
\delta = p - y_{\text{one-hot}}, \qquad \text{onde } y_{\text{one-hot}} = \big(\mathbf{1}[1=y],\dots,\mathbf{1}[C=y]\big)^\top
$$

Ou seja, $\delta$ é simplesmente o vetor de probabilidades previstas menos um vetor com 1 na posição da classe correta e 0 nas demais. Os gradientes da Seção 4, todos de uma vez, ficam:
$$
\frac{\partial L}{\partial W} = \delta\, x^\top \in \R^{C\times d}, \qquad
\frac{\partial L}{\partial b} = \delta \in \R^C
$$

\begin{sanidade}
$\delta \in \R^{C\times 1}$ e $x^\top \in \R^{1\times d}$: o produto $\delta x^\top$ é um produto externo (\emph{outer product}), resultando numa matriz $C\times d$ --- não um escalar (isso seria $\delta^\top x$, que nem faz sentido dimensional aqui) e não um vetor. A entrada $(k,i)$ dessa matriz é $\delta_k \cdot x_i$, que é exatamente $\partial L/\partial W_{k,i}$: o quanto a perda muda se eu mexer no peso que conecta a feature $i$ à classe $k$.

Sempre que calcular um gradiente, cheque se a forma bate com a forma do parâmetro original: $\partial L/\partial W$ tem que ter a mesma forma de $W$ ($C\times d$), assim como $\partial L/\partial b$ tem que ter a mesma forma de $b$ ($C$). Se as dimensões não baterem, há um erro na conta.
\end{sanidade}

\section{De volta ao dataset e ao mini-batch}

Como $L$ (ou $L_{\mathcal{B}}$) é uma média de perdas individuais, seu gradiente é a média dos gradientes individuais --- a derivada de uma soma é a soma das derivadas:
$$
\frac{\partial L_{\mathcal{B}}}{\partial W} = \frac{1}{B}\sum_{i\in\mathcal{B}} \delta^{(i)} \big(x^{(i)}\big)^\top, \qquad
\frac{\partial L_{\mathcal{B}}}{\partial b} = \frac{1}{B}\sum_{i\in\mathcal{B}} \delta^{(i)}
$$

\subsection*{A forma que aparece no código}

Empilhando os exemplos do mini-batch em matrizes --- $X_{\mathcal{B}} \in \R^{B\times d}$ (uma linha $\big(x^{(i)}\big)^\top$ por exemplo) e $\Delta \in \R^{B\times C}$ (uma linha $\big(\delta^{(i)}\big)^\top$ por exemplo) --- a soma acima vira um único produto de matrizes:
$$
\frac{\partial L_{\mathcal{B}}}{\partial W} = \frac{1}{B}\Delta^\top X_{\mathcal{B}} \in \R^{C\times d}, \qquad
\frac{\partial L_{\mathcal{B}}}{\partial b} = \frac{1}{B}\sum_{i\in\mathcal{B}} \Delta_{i,:} \in \R^C
$$

\begin{revisao}
Assim como em $W$ (Seção 1), $X_{\mathcal{B}}$ empilha vetores coluna ($x^{(i)}$) como \emph{linhas}, o que exige transpô-los: a linha $i$ de $X_{\mathcal{B}}$ é $\big(x^{(i)}\big)^\top$, não $x^{(i)}$. A definição de $\Delta$ já trazia essa transposição explícita; agora as duas ficam no mesmo padrão.
\end{revisao}

\begin{intuicao}
É essa vetorização --- trocar um laço \texttt{for} sobre $B$ exemplos por um único produto de matrizes --- que faz frameworks como PyTorch e TensorFlow rodarem rápido em GPU. Mas por trás de cada chamada a \texttt{loss.backward()} está exatamente a conta que fizemos à mão nas Seções 4 e 5: nada mágico, só álgebra linear organizada.
\end{intuicao}

\section*{Para praticar}
\begin{enumerate}[leftmargin=1.3cm]
  \item Refaça, com suas próprias palavras (sem olhar a caixa \emph{Para os curiosos}), a derivação de $\partial L/\partial z_k = p_k - \mathbf{1}[k=y]$ a partir de $L = -z_y + \log(S)$. É um ótimo treino de regra da cadeia.
  \item Para $C=3$ classes e um exemplo com $z=(2.0, 1.0, 0.1)$ e classe correta $y=1$ (primeira classe), calcule $p$, $L$ e $\delta$ na mão.
  \item Implemente as fórmulas vetorizadas da Seção 6 em NumPy, para um mini-batch de $B=4$ exemplos e $C=3$ classes, e confira as formas (\texttt{.shape}) de cada gradiente contra a checagem de sanidade da Seção 5.
\end{enumerate}

\end{document}
